\section{Strip-confined models: compactness and controlled negative couplings}\label{sec:strips}

The Lee--Yang real-zero class has strong closure properties; see \cite{LY52,NW19}. Merely replacing real zeros by a common horizontal strip does not preserve those properties. The counterexample and repair in \cite{R1002} are therefore part of the mathematical result, not optional cautions about numerical fitting.

\subsection{A moment-convergence counterexample}
For $a\ge0,w>0$,
\[
 f_{a,w}(z)=\frac{a+\cos(wz)}{a+1}
\]
is the characteristic function of a symmetric three-point law. If $0\le a\le1$ all its zeros are real. If $a>1$, the zeros are
\[
 \frac{(2j+1)\pi\pm\ii\arcosh(a)}{w},\qquad j\in\Z.
\]
Fix $b>0$, set $a_n=\cosh(bn)$ and $v_n=n^2/(1+a_n)$, and, for large $n$, define
\begin{equation}\label{eq:rare-jump}
 f_n(z)=\cos(\sqrt{1-v_n}\,z)\,
          \frac{a_n+\cos(nz)}{1+a_n}.
\end{equation}
These are symmetric characteristic functions of variance exactly one. All their zeros lie in $\abs{\Im z}\le b$. Every fixed moment converges to the corresponding moment of the fair sign law: the rare-jump contribution to the $2k$th moment is $n^{2k}/(1+\cosh(bn))\to0$, and $v_n\to0$. Yet, for every $y>b$,
\[
 f_n(\ii y)\ge\frac{\cosh(bn)+\cosh(ny)}{1+\cosh(bn)}\longrightarrow\infty.
\]
Thus even convergence of every fixed moment does not give local boundedness in the region where zero exclusion would be transferred.

\begin{theorem}[One exponential budget suffices]\label{thm:strip-transfer}
Let $\nu$ be the actual theta probability law with characteristic function $F_\theta(z)=\Xi(z)/\Xi(0)$. Let $\nu_m$ be symmetric probability measures with entire characteristic functions $f_m$. Suppose, for fixed $0\le b<1/2<R$ and $C<\infty$, that all zeros of $f_m$ lie in $\abs{\Im z}\le b$, that
\begin{equation}\label{eq:one-exp-budget}
 \int e^{R\abs x}\dd\nu_m(x)\le C,
\end{equation}
and that every fixed moment converges to the corresponding theta moment. Then $\Xi$ has no zeros with $\abs{\Im z}>b$. Hence ordinary zeta has no nontrivial zero with $\Re\rho>1/2+b$.
\end{theorem}
\begin{proof}
For $\abs{\Im z}<R$, $\abs{f_m(z)}\le C$, so the functions form a normal family on this connected strip. Every subsequential limit has at zero the prescribed derivatives, since local uniform convergence permits differentiation. It therefore equals $F_\theta$ near zero, and then on the strip by analytic uniqueness. This identifies every subsequential limit and proves local uniform convergence of the sequence. Hurwitz applies on each of $b<\Im z<R$ and $-R<\Im z<-b$. The limit is not identically zero because its value at zero is one. All possible nontrivial zeta zeros already correspond to $\abs{\Im z}<1/2<R$, which proves the conclusion at every height.
\end{proof}

It suffices to certify the zero exclusion on expanding horizontal windows inside $\abs{\Im z}<R$, because each compact eventually lies in such a window. If the first $L$ moments match exactly and both measures satisfy the exponential budget $C$, the elementary Taylor-tail estimate is
\[
 \sup_{\abs z\le r}\abs{f_m(z)-F_\theta(z)}
 \le2C(r/R)^{L+1},\qquad r<R.
\]
Moment mismatches add their explicit finite Taylor sum. This is a disc estimate, not a uniform assertion on an unbounded horizontal strip.

The finite three-point factors above describe a test class and the counterexample, not a successful cofinal theta construction. The separate product obstruction in \cite[strip proof]{R1001} further warns against treating products of such elementary laws as a universal approximation class. The correlated graph construction below has different hypotheses and is not excluded by that product argument.

\subsection{A finite-graph strip theorem}
Consider a finite graph without self-loops, with finite real pair couplings $J_e$ and real fields $h_i$. Under the probability law proportional to
\[
 \exp\left(\sum_{e=\{i,j\}}J_e\sigma_i\sigma_j+\sum_i h_i\sigma_i\right),
 \qquad \sigma_i\in\{-1,1\},
\]
let $Y=\sum_iw_i\sigma_i$, $w_i\ge0$, and $\varphi(z)=\E e^{\ii zY}$. Parallel edges may first be combined. The following result is the controlled-negative-coupling extension in \cite{R1003}; its proof gives the needed disk form of the classical contraction argument directly.

\begin{theorem}[Allocatable negative-edge budget]\label{thm:negative-strip}
For each negative edge $e=\{i,j\}$, choose nonnegative endpoint budgets $\alpha_{e,i},\alpha_{e,j}$ satisfying
\begin{equation}\label{eq:edge-budget}
 \tanh\alpha_{e,i}\tanh\alpha_{e,j}\ge\tanh\abs{J_e}.
\end{equation}
Suppose, at each vertex,
\begin{equation}\label{eq:vertex-budget}
 \abs{h_i}+\sum_{e\ni i,\ J_e<0}\alpha_{e,i}\le b w_i.
\end{equation}
Then every zero of the complete characteristic function $\varphi$ lies in $\abs{\Im z}\le b$. There is no restriction on the number of vertices, cycles, or positive couplings.
\end{theorem}
\begin{proof}
For an edge put $c=e^{-2J}>0$. Its two-variable polynomial, up to a nonzero positive constant, is
\[
 p_c(x,y)=1+c(x+y)+xy.
\]
If $c\le1$, it is nonzero when $\abs x,\abs y<1$. For $c<1$, the possible root $y=-(1+cx)/(c+x)$ has modulus greater than one, since
$\abs{1+cx}^2-\abs{c+x}^2=(1-c^2)(1-\abs x^2)>0$; the case $c=1$ factors.

If $c>1$, choose $r_i,r_j\in(0,1)$ with
$1+r_ir_j\ge c(r_i+r_j)$. This implies $r_i,r_j<1/c$. Minimization on a circle gives, for $\abs x\le r_i$,
\[
 \left|\frac{1+cx}{c+x}\right|
 \ge\frac{1-c\abs x}{c-\abs x}
 \ge\frac{1-cr_i}{c-r_i}\ge r_j.
\]
Thus the edge polynomial is nonzero in the open bidisk. With $r_i=e^{-2\alpha_i}$ and $r_j=e^{-2\alpha_j}$, the radius condition is equivalent to \eqref{eq:edge-budget}, because
\[
 \frac{1+r_ir_j}{r_i+r_j}
 =\frac{1+\tanh\alpha_i\tanh\alpha_j}
        {1-\tanh\alpha_i\tanh\alpha_j}.
\]

The disk contraction step is elementary. If $P(x,y)=A+Bx+Cy+Dxy$ is nonzero on $\abs x<r_1,\abs y<r_2$, then $A+Dz$ is nonzero for $\abs z<r_1r_2$. Indeed $A\ne0$, and the roots of $P(r_1t,r_2t)$ have modulus at least one. Their product gives $\abs D r_1r_2\le\abs A$; if $D=0$ the conclusion is immediate. Other variables can be held fixed throughout, so the contraction can be iterated.

Assign separate variables to edge incidences, multiply all edge polynomials, and contract the incidences at each vertex. An isolated vertex supplies $1+z_i$. The result is exactly the graph subset polynomial
\[
 \mathcal Z(\boldsymbol z)
 =\sum_{S\subseteq V}
  \left(\prod_{e:\,\abs{e\cap S}=1}e^{-2J_e}\right)\prod_{i\in S}z_i.
\]
Every physical spin configuration remains. The contraction removes only inconsistent incidence assignments introduced in the factorization. It proves nonvanishing when all $\abs{z_i}<R_i$, where
$R_i=\exp(-2\sum_{e\ni i,J_e<0}\alpha_{e,i})$. Spin reversal gives nonvanishing also when all $\abs{z_i}>1/R_i$.

Set $z_i=e^{2h_i+2\ii zw_i}$. For positive weights, \eqref{eq:vertex-budget} puts these activities strictly inside the respective disks when $\Im z>b$, and strictly outside the reciprocal disks when $\Im z<-b$. The additional exponential and normalizing factors are nonzero. A zero-weight vertex has zero field and no incident negative edge. Replacing those zero weights by positive $\eps$ and then letting $\eps\downarrow0$ gives the result by Hurwitz; the finite entire limit has value one at zero and is not identically zero.
\end{proof}

A simple symmetric allocation is
\[
 \alpha_{e,i}=\alpha_{e,j}
 =\tfrac12\arcosh(e^{2\abs{J_e}}),
\]
but asymmetric allocations are allowed. Thus a larger observable weight can pay a larger share of an incident negative-edge cost. This theorem does not assign the ordinary Lee--Yang real-zero property to an antiferromagnet; it pays an explicit nonzero strip width.

\subsection{A reference-variance bound and its continuous barrier}
Let the positive reference model have zero fields and couplings $J_e^+=\max(J_e,0)$. Put
\[
 C_{ij}=\E_+(\sigma_i\sigma_j),\qquad V_+=w^{\mathsf T}Cw,
 \qquad K_-=\sum_{J_e<0}\abs{J_e}.
\]
The correlation matrix is positive semidefinite, has diagonal one, and has nonnegative entries. The zero-field ferromagnetic Gaussian estimate is
\begin{equation}\label{eq:ferro-Gaussian}
 \E_+e^{\sum_i q_i\sigma_i}\le e^{q^{\mathsf T}Cq/2},\qquad q_i\ge0.
\end{equation}
This is the classical Lee--Yang consequence used in \cite[strip theorem]{R1001}. For example, apply the real-zero product to the finite weighted observable $\sum_iq_i\sigma_i$ and use $\log(1+x)\le x$ on its moment generating product. Arbitrary field signs can be replaced by absolute values on the right, by the nonnegative zero-field spin correlations.

\begin{proposition}[Reference variance and entire growth]\label{prop:reference-variance}
Under \eqref{eq:edge-budget}--\eqref{eq:vertex-budget}, let $v=\E Y^2$. Then
\begin{align}
 v&\ge V_+e^{-b^2V_+},\label{eq:variance-barrier}\\
 \abs{\varphi(z)}&\le
 \exp\left[V_+\left(b^2+b\abs{\Im z}+\tfrac12\abs{\Im z}^2\right)\right].
 \label{eq:graph-envelope}
\end{align}
Here $v$ is a second moment; no centering is needed for these inequalities.
\end{proposition}
\begin{proof}
Write $d_i=\sum_{e\ni i,J_e<0}\alpha_{e,i}$ and $q_i=\abs{h_i}$. The edge condition implies
\[
 \abs{J_e}\le\tfrac12\{
 \log\cosh(\alpha_{e,i}+\alpha_{e,j})
 -\log\cosh(\alpha_{e,i}-\alpha_{e,j})\}
 \le\alpha_{e,i}\alpha_{e,j},
\]
where the last bound integrates $\tanh t\le t$. Hence
\[
 2K_-\le\sum_i d_i^2\le d^{\mathsf T}Cd,
 \qquad 2K_-+\tfrac12q^{\mathsf T}Cq\le b^2V_+.
\]
For the second inequality use $d+q\le bw$ componentwise and nonnegative entries of $C$.

Let $A=\sum_{J_e<0}J_e\sigma_i\sigma_j$ and $H=\sum_i h_i\sigma_i$. Global spin reversal of the reference law gives
$\E_+[Y^2e^{A+H}]=\E_+[Y^2e^A\cosh H]\ge e^{-K_-}V_+$. Its normalizer is at most $e^{K_-+q^{\mathsf T}Cq/2}$ by \eqref{eq:ferro-Gaussian}. Division proves \eqref{eq:variance-barrier}.

The normalizer is also at least $e^{-K_-}\E_+e^H\ge e^{-K_-}$. With $y=\abs{\Im z}$, the absolute numerator is at most
$e^{K_-}\exp((q+yw)^{\mathsf T}C(q+yw)/2)$. Use the previous budget and $q\le bw$ to obtain \eqref{eq:graph-envelope}.
\end{proof}

A continuous deformation preserving $v$, the band budget $b$, and continuity of $V_+$ cannot escape its small-reference-variance component when it begins at a zero-field ferromagnet. If $b^2v<1/e$, the inequality $v\ge V_+e^{-b^2V_+}$ excludes the open interval between the two roots of $Ve^{-b^2V}=v$. Continuity prevents crossing that interval.

For the physical theta normalization, take $b=1/3$ and use $v_\theta<4621/100000$. The explicit barrier $r=2323/50000=0.04646$ satisfies
\[
 r-r^2/9-4621/100000=\frac{228671}{22500000000}>0.
\]
Since $e^{-x}\ge1-x$, a first crossing $V_+=r$ would contradict \eqref{eq:variance-barrier}. Thus $V_+<r$ throughout such a continuous component. Graph changes are allowed when they preserve this continuity, for example the insertion of zero edges or zero-weight vertices.

A cofinal sequence of \emph{symmetric}, theta-matching models in this component would have the common entire envelope \eqref{eq:graph-envelope}. Montel, coefficient identification, and Hurwitz would then give the ordinary $5/6$ boundary. No such cofinal sequence, or negative-coupling theta-matching finite model, is supplied by \cite{R1003}. The contribution is an enlarged admissible class with a paid strip and compactness budget, not an achieved zero-free region.
